%% toc: P2 17 17.4,17.5,17.6
\section{First-order Hilbert deduction, Kleene's predicate calculus and tableau equality}
\label{sec:p02-w02-p02-025-fol-calculi}

This section completes three pieces of first-order proof theory that the
surrounding treatment uses but does not write out.  First, it gives the two
propositional derivations on which the Hilbert-FOL deduction theorem
(\cref{thm:p02-hilbert-fol-deduction}) rests, the identity
\(\varphi\to\varphi\) and the permutation of antecedents, as explicit
axiom-and-modus-ponens sequences with empty premise support, together with
the bookkeeping lemma that lets a premise-independent line be rederived
without the discharged premise.  With these, every step of the deduction
theorem's proof is a written derivation.  Second, it presents Kleene's own
predicate calculus, whose quantifier postulates 9--12 and deduction theorem
differ in form from the registered sheet, and proves that deduction theorem
under Kleene's condition that free variables be held constant.  Third, it
records where the equality rules of first-order tableaux, their
side conditions and their adequacy proofs are found, and adds those rules to
the registered tableau sheet.

\emph{Inputs.}  For the Hilbert results: a set \(\Gamma\) of formulas of the
equality-free fragment of \texttt{sheet:\allowbreak hilbert-fol}, whose
primitives are \(\neg\), \(\to\), \(\forall\) and \(\exists\), a formula
\(\varphi\), and a
dependency-annotated derivation in the sense of
\cref{def:p02-hilbert-fol-dependencies}.  For Kleene's calculus: a list of
assumption formulas of his predicate calculus and a deduction from it.  For
tableaux: a finite set of first-order sentences, with or without equality.
\emph{Conclusions.}  Explicit derivations of \(\varphi\to\psi\) from
\(\Gamma\) with support not containing \(\varphi\); Kleene's deduction
theorem; and the location and rule form of the equality tableau calculus with
its soundness, countermodel and completeness results.

\subsection{Identity, permutation and supports}
\label{sec:p02-w02-p02-025-kernels}

Throughout this subsection \(\mathsf H_{\mathsf{FOL}}\) is the calculus of
\texttt{sheet:\allowbreak hilbert-fol} with the dependency annotation of
\cref{def:p02-hilbert-fol-dependencies}: each line is a pair \((\chi,D)\),
where the \emph{support} \(D\) is the finite set of premises on which the line
depends.  An \emph{empty-support derivation} of \(\chi\) is a derivation from
any premise set whose last line is \((\chi,\varnothing)\); it uses no premise
line at all, so it is also a derivation from every premise set.  The
derivations in the next two lemmas use only the schemas A1 and A2 and modus
ponens.  They contain no generalization and no existential elimination, so no
variable condition arises, and the metavariables may be replaced by any
fragment formulas, including formulas with free variables or quantifiers.

\begin{hblemma}{thm:p02-w02-p02-025-identity}{proved=annex:p02-w02-p02-025-identity}{Identity with empty support}
For every fragment formula \(\varphi\), the following five lines form an
empty-support derivation of \(\varphi\to\varphi\) in
\(\mathsf H_{\mathsf{FOL}}\):
\[
\begin{array}{lll}
 1. & \varphi\to((\varphi\to\varphi)\to\varphi) & \text{A1}\\
 2. & (\varphi\to((\varphi\to\varphi)\to\varphi))\to
      ((\varphi\to(\varphi\to\varphi))\to(\varphi\to\varphi)) & \text{A2}\\
 3. & (\varphi\to(\varphi\to\varphi))\to(\varphi\to\varphi) & \text{MP }2,1\\
 4. & \varphi\to(\varphi\to\varphi) & \text{A1}\\
 5. & \varphi\to\varphi & \text{MP }3,4
\end{array}
\]
It uses the same three axiom instances as the identity derivation in
Annex~\ref{annex:p2.calc.2} for the propositional sheet, with \(\varphi\) now
ranging over first-order fragment formulas.
\end{hblemma}

\begin{hblemma}{thm:p02-w02-p02-025-permutation}{proved=annex:p02-w02-p02-025-permutation}{Permutation of antecedents with empty support}
For all fragment formulas \(\varphi,\theta,\psi\) there is an empty-support
derivation in \(\mathsf H_{\mathsf{FOL}}\), using only A1, A2 and modus
ponens, of
\[
 (\varphi\to(\theta\to\psi))\to(\theta\to(\varphi\to\psi)).
\]
Consequently, if a derivation from \(\Gamma\) has a line
\((\varphi\to(\theta\to\psi),E)\), appending that derivation and one modus
ponens step gives a line \((\theta\to(\varphi\to\psi),E)\) with the same
support; and the instance with \(\varphi\) and \(\theta\) interchanged gives
the inverse passage from \(\theta\to(\varphi\to\psi)\) to
\(\varphi\to(\theta\to\psi)\), again with unchanged support.
\end{hblemma}

The derivation has nineteen lines and is written out in the annex.  No
appeal to propositional completeness is involved: the permutation is not
obtained from the cited basis theorem
\cref{thm:p02-calculi-derivability-equivalence} through
\cref{thm:p02-hfa-propositional-kernels}, but from five A1 instances, five A2
instances and nine modus ponens steps.

The deduction theorem's proof also needs two facts about supports.  Given a
derivation \(\langle(\chi_1,D_1),\ldots,(\chi_n,D_n)\rangle\), say that line
\(i\) \emph{uses} line \(j<i\) when line \(i\) is obtained by modus ponens
from line \(j\) and another line, or by generalization or existential
elimination from line \(j\).

\begin{hblemma}{thm:p02-w02-p02-025-support}{proved=annex:p02-w02-p02-025-support}{Concatenation and support restriction}
Let \(\Gamma,\Delta\) be sets of fragment formulas.
\begin{enumerate}
  \item If \(\mathcal D_1\) and \(\mathcal D_2\) are dependency-annotated
        derivations from \(\Gamma\), then the concatenation
        \(\mathcal D_1\mathcal D_2\), with the line references of
        \(\mathcal D_2\) shifted by the length of \(\mathcal D_1\), is a
        dependency-annotated derivation from \(\Gamma\) whose lines carry
        their original supports.
  \item If a dependency-annotated derivation from \(\Gamma\) has a line
        \((\chi,D)\) and \(D\subseteq\Delta\), then there is a
        dependency-annotated derivation from \(\Delta\) ending in
        \((\chi,D)\).  It consists of the lines on which \((\chi,D)\)
        depends through the uses relation, in their original order and with
        their original supports.
\end{enumerate}
\end{hblemma}

Clause~2 is the precise form of the phrase ``its traced subderivation is
already from \(\Gamma\)'' in the proof of the deduction theorem: if the
discharged premise \(\varphi\) is not in the support of a line, the line can
be rederived from the remaining premises, and its generalization or
existential-elimination condition, which mentions only the support, is
unchanged.  With the three lemmas, every step of the deduction theorem is an
explicit derivation.

\begin{hbtheorem}{thm:p02-w02-p02-025-discharge}{proved=annex:p02-w02-p02-025-discharge}{Hilbert-FOL deduction theorem with explicit supports}
Let \(\Gamma\) be a set of fragment formulas and \(\varphi,\psi\) fragment
formulas.
\begin{enumerate}
  \item If a dependency-annotated derivation from \(\Gamma\cup\{\varphi\}\)
        ends in \((\psi,D)\), then a dependency-annotated derivation from
        \(\Gamma\) ends in \((\varphi\to\psi,E)\) for some
        \(E\subseteq D\setminus\{\varphi\}\).
  \item If a dependency-annotated derivation from \(\Gamma\) ends in
        \((\varphi\to\psi,E)\), then one from \(\Gamma\cup\{\varphi\}\) ends
        in \((\psi,E\cup\{\varphi\})\).
\end{enumerate}
In particular
\(\Gamma\cup\{\varphi\}\derives_{\mathsf H_{\mathsf{FOL}}}\psi\) if and only if
\(\Gamma\derives_{\mathsf H_{\mathsf{FOL}}}\varphi\to\psi\), which is
\cref{thm:p02-hilbert-fol-deduction}.  The construction of clause~1 applies
generalization to a line depending on \(\varphi\) only when the original
derivation did, and then the rule's condition already guarantees that the
generalized variable is not free in \(\varphi\).
\end{hbtheorem}

The theorem has the same content as the discharge interface
\cref{thm:p02-hfa-discharge-interface}; the present proof writes each case as
a derivation, using \cref{thm:p02-w02-p02-025-identity} for the discharged
premise, A1 for other premises and axioms, A2 for modus ponens, the
\(\forall2\) schema for generalization, and
\cref{thm:p02-w02-p02-025-permutation} twice for existential elimination.

\begin{hbexample}{def:p02-w02-p02-025-example}{Discharging a premise through generalization}
Let \(P,Q\) be unary relation symbols, \(\Gamma=\{\forall x(Px\to Qx)\}\) and
\(\varphi=\forall x\,Px\).  A derivation of \(\forall x\,Qx\) from
\(\Gamma\cup\{\varphi\}\), with supports, is:
\[
\begin{array}{llll}
 1. & \forall x(Px\to Qx) & \{\forall x(Px\to Qx)\} & \text{premise}\\
 2. & \forall x(Px\to Qx)\to(Px\to Qx) & \varnothing & \forall1\\
 3. & Px\to Qx & \{\forall x(Px\to Qx)\} & \text{MP }2,1\\
 4. & \forall x\,Px & \{\forall x\,Px\} & \text{premise}\\
 5. & \forall x\,Px\to Px & \varnothing & \forall1\\
 6. & Px & \{\forall x\,Px\} & \text{MP }5,4\\
 7. & Qx & \{\forall x(Px\to Qx),\forall x\,Px\} & \text{MP }3,6\\
 8. & \forall x\,Qx & \{\forall x(Px\to Qx),\forall x\,Px\} & \text{Gen }7
\end{array}
\]
Gen at line~8 is permitted because \(x\) is free in neither premise of the
support.  The construction of \cref{thm:p02-w02-p02-025-discharge} turns line
\(i\) into a line \(\varphi\to\chi_i\).  Lines 1, 2 and 5, a premise other
than \(\varphi\) and two axioms, become by A1 and modus ponens
\(\varphi\to\chi_i\) with the same support.  Line~4 is \(\varphi\) and becomes
the identity derivation of \(\forall x\,Px\to\forall x\,Px\), with empty
support.  Line~3 comes by modus ponens, and A2 turns it into
\(\forall x\,Px\to(Px\to Qx)\) with support \(\{\forall x(Px\to Qx)\}\).
Lines 6 and 7 come by modus ponens as well, and A2 turns them into
\(\forall x\,Px\to Px\) with empty support and
\(\forall x\,Px\to Qx\) with support \(\{\forall x(Px\to Qx)\}\).  At line~8
the support contains \(\varphi\), so the construction generalizes
\(\forall x\,Px\to Qx\), which is permitted because \(x\) is free in no
member of \(\{\forall x(Px\to Qx)\}\), and applies the \(\forall2\) instance
\(\forall x(\forall x\,Px\to Qx)\to(\forall x\,Px\to\forall x\,Qx)\), whose
side condition holds because \(x\notin\FV(\forall x\,Px)\).  The result is
\(\Gamma\derives_{\mathsf H_{\mathsf{FOL}}}\forall x\,Px\to\forall x\,Qx\)
with support \(\{\forall x(Px\to Qx)\}\).
\end{hbexample}

\begin{hbwarning}{def:p02-w02-p02-025-unrestricted}{The deduction theorem fails for unrestricted generalization}
Suppose generalization were allowed on every line.  From the premise \(Px\),
one step would give \(\forall x\,Px\), so
\(\{Px\}\) would derive \(\forall x\,Px\).  The formula
\(Px\to\forall x\,Px\) is not valid: in a structure with domain \(\{0,1\}\)
and \(P\) interpreted as \(\{0\}\), the assignment sending \(x\) to \(0\)
makes \(Px\) true and \(\forall x\,Px\) false.  By the soundness theorem
\cref{thm:p02-hfa-soundness} it is not derivable from the empty set.  So the
equivalence of \cref{thm:p02-w02-p02-025-discharge} would fail.  On the
registered sheet the step is blocked: the line \((Px,\{Px\})\) has \(x\) free
in a member of its support.  Kleene's form of the same restriction is that
free variables be held constant for the discharged assumption
(\cref{def:p02-w02-p02-025-kleene-dependence}).
\end{hbwarning}

\subsection{Kleene's predicate calculus and his deduction theorem}
\label{sec:p02-w02-p02-025-kleene}

The registered sheet takes Mendelson's universal base and adds an
existential schema and rule; it is not literally Kleene's system.  Kleene's
own postulates for the predicate calculus, Group~A of his formal system, are
given in \cite[\S 19, pp.~81--82]{Kleene1952}, and his deduction theorem for
the predicate calculus in \cite[\S\S 21--22, pp.~90--98]{Kleene1952}.  Page
numbers refer to the held 1974 seventh reprint.  The presentation below keeps
Kleene's notation: \(\supset\) for implication, \(\mathbin{\&}\) for
conjunction, \(\lor\), \(\neg\), \(\forall\), \(\exists\), with
\(\mathbin{\&}\) and \(\lor\) binding more tightly than \(\supset\), as in his
displays; \(A(x)\) and \(A(t)\) denote a formula and the result of
substituting \(t\) for the free occurrences of \(x\) in it.  In this
subsection \(\derives\) is Kleene's deducibility relation.

\begin{hbdefinition}{def:p02-w02-p02-025-kleene-postulates}{Kleene's postulates for the predicate calculus}
In postulates 1--8, \(A\), \(B\) and \(C\) are formulas.  In postulates
9--12, \(x\) is a variable, \(A(x)\) is a formula, \(C\) is a formula that does
not contain \(x\) free, and \(t\) is a term that is free for \(x\) in
\(A(x)\) \cite[\S 19, p.~81]{Kleene1952}.

\emph{Group A1 (propositional calculus)} \cite[\S 19, p.~82]{Kleene1952}:
\[
\begin{array}{ll}
 \text{1a.} & A\supset(B\supset A).\\
 \text{1b.} & (A\supset B)\supset((A\supset(B\supset C))\supset(A\supset C)).\\
 \text{2.}  & \text{Rule: from } A \text{ and } A\supset B \text{, infer } B.\\
 \text{3.}  & A\supset(B\supset A\mathbin{\&}B).\\
 \text{4a.} & A\mathbin{\&}B\supset A.\\
 \text{4b.} & A\mathbin{\&}B\supset B.\\
 \text{5a.} & A\supset A\lor B.\\
 \text{5b.} & B\supset A\lor B.\\
 \text{6.}  & (A\supset C)\supset((B\supset C)\supset(A\lor B\supset C)).\\
 \text{7.}  & (A\supset B)\supset((A\supset\neg B)\supset\neg A).\\
 \text{8}^{\circ}. & \neg\neg A\supset A.
\end{array}
\]
\emph{Group A2, additional postulates for the predicate calculus,}
\cite[\S 19, p.~82]{Kleene1952}:
\[
\begin{array}{ll}
 \text{9.}  & \text{Rule: from } C\supset A(x) \text{, infer } C\supset\forall xA(x).\\
 \text{10.} & \forall xA(x)\supset A(t).\\
 \text{11.} & A(t)\supset\exists xA(x).\\
 \text{12.} & \text{Rule: from } A(x)\supset C \text{, infer } \exists xA(x)\supset C.
\end{array}
\]
The \emph{axioms} of the predicate calculus are the formulas of the forms 1a,
1b, 3--8, 10 and 11; postulates 2, 9 and 12 are the \emph{rules of
inference}, and a formula is an \emph{immediate consequence} of the premise
or premises of one of them \cite[\S 19, p.~82]{Kleene1952}.  The mark on
postulate~8 is Kleene's; he defers its explanation to his \S 23
\cite[\S 19, p.~82]{Kleene1952}.  Kleene's full number-theoretic system adds postulates 13--21
(Group~B); they are axioms, not rules, and are not used here.
\end{hbdefinition}

\begin{hbdefinition}{def:p02-w02-p02-025-kleene-dependence}{Deduction, dependence, and variables held constant}
Let \(D_1,\ldots,D_l\) be a list of formula occurrences, \(l\geq 0\).
\begin{enumerate}
  \item A \emph{deduction} from the assumption formulas \(D_1,\ldots,D_l\) is
        a finite sequence of one or more formula occurrences, each of which
        is one of \(D_1,\ldots,D_l\), an axiom, or an immediate consequence
        of preceding formulas of the sequence; \(D_1,\ldots,D_l\derives E\)
        means that a deduction from them ends in \(E\)
        \cite[\S 20, p.~87]{Kleene1952}.  An \emph{analysis} records, for each
        member, which of these justifications is used and from which earlier
        members.
  \item Given a deduction \(A_1,\ldots,A_k\) with an analysis,
        \(A_i\) \emph{depends on} the assumption occurrence \(D_j\) when
        \(A_i\) is \(D_j\) in the analysis, or \(A_i\) is an immediate
        consequence of a member that depends on \(D_j\); and only in these
        cases \cite[\S 22, p.~95]{Kleene1952}.
  \item A variable \(y\) is \emph{varied} in the deduction for the assumption
        formula \(D_j\) if \(y\) occurs free in \(D_j\) and the deduction
        applies Rule 9 or Rule 12 with respect to \(y\) (that is, with \(y\) as
        the \(x\) of the rule) to a premise that depends on \(D_j\).
        Otherwise \(y\) is \emph{held constant} for \(D_j\)
        \cite[\S 22, pp.~95--96]{Kleene1952}.
\end{enumerate}
Kleene writes \(\Gamma\) for a finite list of zero or more formula
occurrences and \(\Gamma,A\) for the list with \(A\) appended last.
\end{hbdefinition}

\begin{hbtheorem}{thm:p02-w02-p02-025-kleene-deduction}{proved=annex:p02-w02-p02-025-kleene-deduction}{Kleene's deduction theorem for the predicate calculus}
For the predicate calculus, if \(\Gamma,A\derives B\) with the free variables
held constant for the last assumption formula \(A\), then
\(\Gamma\derives A\supset B\).  Conversely, \(\Gamma\derives A\supset B\)
implies \(\Gamma,A\derives B\).
\end{hbtheorem}

The forward implication is Kleene's Theorem~1 in its concluded form
\cite[\S 22, p.~97]{Kleene1952}; the propositional case, with only Group~A1
in force and no variable condition, is \cite[\S 21, p.~90]{Kleene1952}.  The
proof in the annex follows Kleene's case division: for Rule 9 applied to a
premise depending on \(A\), it passes through \(A\mathbin{\&}C\) and uses the
conjunction postulates 3, 4a and 4b, as Kleene does on page 97, and for
Rule 12 it uses the permutation of antecedents twice, as he indicates on
page 98.  Kleene attributes the first proof of the deduction theorem as a
derived rule to Herbrand \cite[\S 22, p.~98]{Kleene1952}.

\begin{hbwarning}{def:p02-w02-p02-025-kleene-contrast}{Two different quantifier rule forms}
Kleene's Rule 9 has the form ``from \(C\supset A(x)\) infer
\(C\supset\forall xA(x)\)'', with \(x\) not free in \(C\), and places no
condition on the assumption formulas.  The registered sheet's generalization
rule has the form ``from \(\varphi\) infer \(\forall x\,\varphi\)'' and
carries its premise condition in the rule itself, through the support of
\cref{def:p02-hilbert-fol-dependencies}.  So in Kleene's system
\(\derives\) may generalize a variable free in an assumption, and the
restriction appears only as the hypothesis ``held constant'' of the
deduction theorem; on the sheet the restriction is part of
derivability, and the deduction theorem has no extra hypothesis.  Kleene's
1b is \((A\supset B)\supset((A\supset(B\supset C))\supset(A\supset C))\),
not the sheet's A2; Kleene takes \(\mathbin{\&}\), \(\lor\) and \(\exists\) as
primitive with their own postulates, whereas the sheet's language has only
\(\neg,\to,\forall,\exists\).  Neither presentation is a notational variant
of the other, and results proved for one are transferred to the other only
through explicit derivations.
\end{hbwarning}

\subsection{Equality rules for first-order tableaux}
\label{sec:p02-w02-p02-025-tableau-equality}

The first-order tableau calculus of \cref{def:p02-fol-tableau} is
equality-free.  Its extension by the two tableau equality rules, reflexivity
and left-to-right replacement of one occurrence of a closed term, is defined
in \cref{def:p02-w02-p02-014-fol-equality}, with the task bookkeeping of
\cref{def:p02-w02-p02-014-fol-branch}.  That bookkeeping states the
side conditions exactly.  \emph{Freshness}: a delta step uses the parameter
of least index not occurring on the branch, and
\cref{def:p02-w02-p02-014-fol-freshness} shows that dropping this condition
closes a tableau for a satisfiable set.  \emph{Instantiation}: gamma steps and
reflexivity steps range over every closed term of the parameter-extended root
language, not only the terms already on the branch.  \emph{Fairness}: an
open branch counts as saturated only when every available task, including
every reflexivity task and every replacement task, is done, and
\cref{def:p02-w02-p02-014-fol-unfair} shows an unfair open branch that is
not a countermodel.  The rules are displayed with these conditions on
\texttt{sheet:\allowbreak tableaux-fol}.

The adequacy results are proved in
\cref{sec:p02-w02-p02-014-fol-calculus,sec:p02-w02-p02-014-fol-herbrand,sec:p02-w02-p02-014-fol-systematic}
for roots with and without equality: soundness of closure
(\cref{thm:p02-w02-p02-014-fol-soundness}); countermodel extraction from an
open saturated, hence fair, branch through the quotient Herbrand structure
(\cref{thm:p02-w02-p02-014-congruence,thm:p02-w02-p02-014-herbrand-truth,thm:p02-w02-p02-014-fol-extraction});
completeness (\cref{thm:p02-w02-p02-014-fol-completeness}); and the infinite
fair-branch construction by which fair systematic expansion semidecides
unsatisfiability (\cref{thm:p02-w02-p02-014-fol-semidecision}).  The
propositional procedure, by contrast, terminates on every input and decides
satisfiability (\cref{thm:p02-w02-p02-014-termination,thm:p02-w02-p02-014-decision});
\cref{thm:p02-w02-p02-014-fol-no-finite-saturation} shows that the
termination argument fails as soon as a branch carries a universal
occurrence, and \cref{def:p02-w02-p02-014-fol-boundary} states the
decision-versus-semidecision boundary.  This section does not repeat those
proofs.

\subsection{Relationships and sources}
\label{sec:p02-w02-p02-025-relations}

The identity and permutation derivations are the first-order instances of
propositional facts that the propositional Hilbert calculus proves through
its deduction theorem, \cref{thm:p02-hilbert-deduction}; Kleene obtains the
permutation the same way, as two applications of his propositional deduction
theorem \cite[\S 21, p.~93]{Kleene1952}.  The explicit nineteen-line
derivation is written here so that the first-order deduction theorem does
not depend on a lifted kernel whose existence rests on the cited
propositional basis theorem.  The supports of
\cref{def:p02-hilbert-fol-dependencies} are the registered counterpart of
Kleene's dependence relation; Kleene's analysis of which formulas depend on
an assumption occurrence is the source of that bookkeeping.  The universal
part of the registered sheet follows \cite{Mendelson1997}; the existential
schema and rule follow Kleene's postulates 11 and 12
\cite[\S 19, p.~82]{Kleene1952}.  The tableau format is that of
\cite{Smullyan1968}, as recorded on the sheets.
